\subsubsection{对偶空间的例子}\label{subsec:dual-space-examples}

\Needspace{4\baselineskip}
\begin{Example}
\label{ex:lp-conjugate-exponent-dual-functional}
设 $1<p<\infty$ 且 $p^{-1}+q^{-1}=1$. 由 Hölder 不等式, 每个 $f\in L^p(\Omega)$ 都可以通过对应关系
\begin{equation}\label{eq:lp-dual-functional-integral-pairing}
l_f(g):=\int_\Omega\overline{f(\bm{x})}g(\bm{x})\,\mathrm{d}\bm{x}.
\end{equation}
与 $L^q(\Omega)$ 上的一个线性泛函 $l_f$ 相认同. 加入复共轭是为使这一定义与 Hilbert 空间中的内积类似. 显然 $\|l_f\|\leq\|f\|_p$. 现在假设 $f\ne0$, 并设 $f_n$ 为 $\widetilde L^p(\Omega)$ 中的 Cauchy 序列, 在 $L^p$ 中收敛到 $f$. 令 $g_n=f_n|f_n|^{p-2}$. 那么 $g_n\in\widetilde L^q(\Omega)$ 且 $\|g_n\|_q^q=\|f_n\|_p^p$. 进一步得到
\begin{equation}\label{eq:lp-dual-functional-approximating-norm-bound}
|l_f(g_n)-l_{f_n}(g_n)|\leq\|f-f_n\|_p\|g_n\|_q=\|f-f_n\|_p\|f_n\|_p^{p/q}
\end{equation}
以及
\begin{equation}\label{eq:lp-dual-functional-approximating-pairing-value}
l_{f_n}(g_n)=\|f_n\|_p^p.
\end{equation}
由此可得 $l_f(g_n)$ 收敛到 $\|f\|_p^p$, 而由于 $\|g_n\|_q$ 收敛到 $\|f\|_p^{p/q}$, 实际上得到 $\|l_f\|=\|f\|_p$. 因而有从 $L^p(\Omega)$ 到 $L^q(\Omega)$ 对偶空间的等距嵌入. 可以证明, 实际上 $L^q(\Omega)$ 上的所有泛函都可按这种方式表示. 这一定理陈述如下, 不给证明.
\end{Example}

\begin{Theorem}
\label{thm:lp-dual-space-representation}
设 $\Omega$ 为 $\mathbb{R}^m$ 中的局部 Jordan 可测开集, $1<p<\infty$ 且 $p^{-1}+q^{-1}=1$. 那么通过 \eqref{eq:lp-dual-functional-integral-pairing}, $L^p(\Omega)$ 的对偶空间与 $L^q(\Omega)$ 相一致.
\end{Theorem}

\Needspace{4\baselineskip}
\begin{Example}
\label{ex:bounded-continuous-lp-one-dual}
设 $f$ 为 $\Omega$ 上的有界连续函数. 那么有
\begin{equation}\label{eq:bounded-continuous-lp-one-dual-estimate}
\left|\int_\Omega\overline{f(\bm{x})}g(\bm{x})\,\mathrm{d}\bm{x}\right|\leq\|g\|_1\sup_{\bm{x}\in\Omega}|f(\bm{x})|
\end{equation}
对每个 $g\in L^1(\Omega)$ 成立. 此外, 取非负的 $g$, 使其支集位于 $|f|$ 接近其上确界的某点附近, 就可看到
\begin{equation}\label{eq:bounded-continuous-infinity-norm}
\|f\|_\infty:=\sup_{\bm{x}\in\Omega}|f(\bm{x})|
\end{equation}
表示 $f$ 作为 $L^1(\Omega)$ 上线性泛函的范数. 但 $L^1(\Omega)$ 上的线性泛函并非都能按这种方式表示; 例如可以取在某个曲面两侧不连续的 $f$. 作如下定义.
\end{Example}

\begin{Definition}
\label{def:lp-infinity-dual-space}
设 $\Omega$ 为 $\mathbb{R}^m$ 中的局部 Jordan 可测开集. $L^1(\Omega)$ 的对偶空间记为 $L^\infty(\Omega)$.
\end{Definition}

由于试验函数在 $L^1(\Omega)$ 中稠密, 且试验函数意义下的收敛蕴含 $L^1$ 收敛, $L^\infty(\Omega)$ 是一个分布空间. 此外, 由于连续函数的一致极限仍连续, 连续函数不可能在 $L^\infty(\Omega)$ 中稠密; 事实上, $L^\infty(\Omega)$ 不可分. 在测度论中, $L^\infty$ 被刻画为有界可测函数的空间.
